% GENERATED by manuscript/code/build_tables.py -- do not edit by hand.
\begin{table*}[t]
\centering
\small
\caption{Cross-validated performance of a fixed leakage diagnostic (random forest: 300 trees, unrestricted depth) under three fold-assignment schemes in the construction sample.}
\label{tab:leakage}
\begin{tabular}{lrrrr}
\toprule
Fold assignment & Brier score & AUROC & Log-loss & Folds \\
\midrule
Random-event (naive) & 0.2048 & 0.6952 & 0.9640 & 5 \\
Event-grouped (control) & 0.2009 & 0.7058 & 0.9223 & 5 \\
Sequence-grouped (primary) & 0.1990 & 0.7076 & 0.9709 & 5 \\
\bottomrule
\end{tabular}
\par\vspace{2pt}\footnotesize\emph{Note.} The diagnostic is distinct from the tuned temporal-holdout baseline B4 (depth 3). Using unrounded values, the sequence-grouped minus random-event Brier difference is -0.0059; the random-event minus sequence-grouped AUROC difference is -0.0124. The preregistration required uncertainty but did not define the paired interval estimator; the executed report retained point estimates only. This deviation is disclosed, and no post-hoc interval or equivalence claim is added. Because there is one row per event, event grouping imposes the same grouping constraint as random-event assignment, but its deterministic fold allocation is different and can produce different scores.
\end{table*}
